\subsection{A further class-two lift with the complete homology action retained}

\begin{statement}{D.51}{Proved}{The moving class-two kernel lifts and the invariant residual is exact}
Keep the actual D.49 groups and the chosen complete D.50
automorphism, including its two conjugators. Put
\[
 \begin{gathered}
 n=|G_3|,\qquad q=|A|,\qquad m={n\over6q},\qquad
 \rho=1+mq,\qquad w=m(q-1),\\
 \mathfrak R_{50}=\mathfrak R_h\cap\mathfrak K_\infty,
 \qquad V=U/2U,\qquad W=U^-/2U^- .
 \end{gathered}
 \tag{D.51.1}
\]
Here \(q\) is the order of the same actual nonzero
elementary three-subgroup \(A\); no numerical choice
of that arithmetic subgroup is substituted.
Let \(\mathcal T_2\) be D.50's actual class-two,
exponent-four quotient, with kernel
\(\mathcal D\) over \(G_3\). Thus
\[
 \mathcal D^{\rm ab}=(\mathbb Z/4)^\rho,\qquad
 \mathcal D'=\bigwedge\nolimits^2_{\mathbb F_2}V .
\]
The exact image, strengthening D.50's inclusion, is
\[
 \operatorname{im}(\mathfrak R_{50}\to\mathcal T_2)
       =Z:=W\wedge V .
 \tag{D.51.2}
\]
Conjugation on \(\mathcal D'\) factors through \(G_3\).
The normal odd subgroup \(A\) therefore gives
canonical \(G_3\)-submodules
\[
 \begin{gathered}
 Z_0=e_A Z=Z^A,\qquad Z_-=(1-e_A)Z,\qquad Z=Z_0\oplus Z_-,\\
 \dim Z_0={m^2(q-1)\over2}>0,\qquad
 \dim Z_-={w\rho\over2}>0 .
 \end{gathered}
 \tag{D.51.3}
\]
All dimensions here are over \(\mathbb F_2\).

Define the actual further marked quotient and its source kernel by
\[
 \mathcal T_-=\mathcal T_2/Z_0,\qquad
 \mathfrak K_-=\ker(\widehat B\to\mathcal T_-).
 \tag{D.51.4}
\]
The maximal common marked quotient of the D.50
source \(\mathcal Z_h\) and \(\mathcal T_-\) is exactly
\(\mathcal T_2/Z\). The fixed D.50 action and its
complete conjugator pair extend to a peripheral
automorphism of the strictly larger source
\[
 \mathcal Y_h
   :=\widehat B/(\mathfrak R_{50}\cap\mathfrak K_-)
    =\mathcal Z_h
       \times_{\mathcal T_2/Z}\mathcal T_- .
 \tag{D.51.5}
\]
Its multiplier is one and it fixes \(t=sr\).
It preserves every D.50 homology level and the entire
retained D.49 action, parameter, CM pair, full
metabelian action and global-sign congruence action.
The new projection and full-source kernels are
\[
 \begin{gathered}
 \ker(\mathcal Y_h\to\mathcal Z_h)=Z_-,
 \qquad
 \ker(\mathcal Y_h\to\mathcal T_-)
       =\mathfrak K_-\mathfrak R_{50}/\mathfrak R_{50},\\
 \mathfrak R_{51}:=\ker(\widehat B\to\mathcal Y_h)
       =\mathfrak R_{50}\cap\mathfrak K_-,
 \qquad
 \operatorname{im}(\mathfrak R_{51}\to\mathcal T_2)=Z_0 .
 \end{gathered}
 \tag{D.51.6}
\]
In particular the next remaining image is exact and
nonzero. The orders of the two finite comparison
quotients are
\[
 \begin{aligned}
 |\mathcal T_2/Z|
     &=n\,2^{\,2\rho+\binom{m+1}{2}},\\
 |\mathcal T_-|
     &=n\,2^{\,2\rho+\binom{\rho}{2}-m^2(q-1)/2}.
 \end{aligned}
 \tag{D.51.7}
\]
The lift uses actual central corrections in \(Z_-\).
The induced exterior action preserves \(Z_0\), so the
map factors through \(\mathcal T_-\); its unchanged
invertible Frattini action proves generation.
These two points are included in the proof.

There are two separate runnable finite calculations.
They retain the exact order-\(1080\) D.50 certificate,
with \(\rho=181\), \(q=3\), \(m=60\), and \(w=120\).
The class-two lift has exact dimensions
\[
            \dim Z=14460,\qquad
            \dim Z_0=3600,\qquad
            \dim Z_-=10860 .
 \tag{D.51.8}
\]
All \(181\) actual free-generator words and the
unchanged homology matrix are checked. The corrected
cusp word has \(4\,159\,444\) letters in \(B\);
direct collection verifies its equality to an
element of \(Z_0\), hence identity in \(\mathcal T_-\).
Both exact torsion relations are also checked.

The second calculation exhibits the reason for
retaining the invariant residual. For this particular
finite D.50 homology action, the ribbon intersection
form has rank \(146\), but its preservation defect
has rank \(72\), with a nonzero entry at \((0,6)\).
Consequently that specified finite homology action
has no multiplier-one peripheral lift to the
unquotiented class-two source. This is an exclusion
for the separate order-\(1080\) certificate. It does
not evaluate or exclude the actual arithmetic
\(G_3\) action in (D.51.1).
Full profinite GT0 realization remains undecided.

The manuscript passages are Part D(e), section 51.16,
folio 734 verso, scan 174, and the two-torus/Fermat
follow-ups in section 51.15, folios 725--727,
scans 156--160.
\end{statement}
\begin{proof}
\emph{The exact class-two residual.}
Write \(D=\mathcal D\) in this paragraph. Its abelianization
is the free \(\mathbb Z/4\)-module \(U/4U\), and the
alternating commutator pairing is the stated
identification \(D'=\bigwedge^2 V\).
Let \(H\) be the image of \(C_h\) in \(\mathcal T_2\).
It is normal, projects onto \(A\), and has a finite
two-group kernel. Choose an odd-order complement
\(\widetilde A\) in \(H\). Its conjugation action on
\(D^{\rm ab}\) and \(D'\) is the original \(A\)-action:
different lifts differ by \(D\), which acts trivially
on those two modules.

Put \(E=[D,\widetilde A]\). It is normal in \(D\)
and stable under \(\widetilde A\). For example,
the product identity for \([d_1d_2,a]\) expresses
the \(d_1\)-conjugate of \([d_2,a]\) as a product
of two such commutators. Normality of \(H\) gives
\(E\subset H\). Its abelian image is
\[
       \operatorname{im}(E\to U/4U)
           =(1-e_A)(U/4U)=U^-/4U^- .
 \tag{D.51.9}
\]
Commuting this image with all of \(D\) proves
\(Z=W\wedge V\subset E\cap D'\).

We prove the reverse inclusion by an actual quotient map.
Let \(D_+\) be the free class-two exponent-four group
on \(\rho-w\) generators. Choose a basis of
\(e_A(U/4U)\). The projection
\[
                p:U/4U\longrightarrow e_A(U/4U)
\]
lifts to a homomorphism \(D\to D_+\): choose images
of the free generators lifting its basis coordinates.
The homomorphisms with this prescribed abelian map
form an affine space for
\(\operatorname{Hom}_{\mathbb F_2}(V,D_+')\).
Indeed the images of each generator can differ
independently by a central element of exponent two.
Precomposition by \(\widetilde A\) preserves this
affine space, because \(p\) is \(A\)-invariant.
Averaging any affine point over the odd group
\(\widetilde A\) gives a fixed point. This produces
an actual \(A\)-invariant surjection
\[
                      \pi:D\longrightarrow D_+ .
 \tag{D.51.10}
\]
Surjectivity follows from its onto Frattini map.
It kills \(E\), and on commutators its map is
\(\bigwedge^2(V\to V^A)\), whose kernel is exactly
\(W\wedge V\). Hence
\[
                            E\cap D'=Z .
 \tag{D.51.11}
\]

For completeness, the needed coprime commutator
identity in this group follows directly from
(D.51.9)--(D.51.11). Put \(E_1=[E,\widetilde A]\).
Its abelian image is still \(U^-/4U^-\), by odd
averaging. Its own commutators therefore contain
\(\bigwedge^2W\). The other summand
\(W\wedge V^A\) of \(Z\) has zero \(A\)-invariants,
so it equals its commutators with \(A\) and lies
in \(E_1\). Thus \(Z\subset E_1\), and the equality
of abelian images gives \(E_1=E\).
The quotient \(H/E\) is the product of its
two-group kernel and its odd complement, since
that complement centralizes \(D/E\). It is
nilpotent. Every nilpotent quotient of \(H\)
kills \([E,\widetilde A]=E\), because its two-
and three-parts commute. Thus \(E\) is the exact
pronilpotent residual of \(H\). D.49's proved
residual-image comparison gives
\[
              \operatorname{im}(\mathfrak R_h
                                      \to\mathcal T_2)=E .
 \tag{D.51.12}
\]

Now take the actual intersection with
\(\mathfrak K_\infty\). Its image lies in
\(E\cap D'=Z\). Conversely, take \(z\in Z\)
and a preimage \(b\in\mathfrak R_h\).
Its image in \(U\) lies both in \(U^-\) and in
\(4U\), hence in \(4U^-\). By D.50, choose
\(c\in\mathfrak R_h\) whose image in \(U\)
is one quarter of that image. The element
\(bc^{-4}\) then lies in
\(\mathfrak R_h\cap\mathfrak K_\infty\).
It still maps to \(z\), because \(D\) has
exponent four. This proves (D.51.2) using the
actual source preimages.

\emph{The invariant and moving dimensions.}
The actual graph calculation of D.49 gives, as
an \(A\)-module,
\[
              V\simeq\mathbb F_2\oplus\mathbb F_2[A]^m .
 \tag{D.51.13}
\]
Consequently \(V^A\) has dimension \(m+1\), and
each nontrivial character of the elementary
three-group \(A\) has multiplicity \(m\) after
extension to \(\mathbb F_4\). The nontrivial
characters occur in distinct inverse pairs.
In \(\bigwedge^2W\), each inverse pair contributes
exactly \(m^2\) invariant dimensions. The
summand \(W\wedge V^A\) has no invariants.
Extension of the coefficient field preserves
dimensions, so
\[
 \begin{aligned}
 \dim Z_0&={q-1\over2}\,m^2,\\
 \dim Z_-&=w(m+1)+\binom w2-{q-1\over2}\,m^2
             ={w\rho\over2}.
 \end{aligned}
 \tag{D.51.14}
\]
The central idempotent \(e_A\) commutes with
the \(G_3\)-action, because \(A\) is normal.
Thus both subspaces are actual normal subgroups
of \(\mathcal T_2\). They are nonzero since
\(q\ge3\) and \(m>0\). The collected order
of \(D\) is \(2^{2\rho+\binom\rho2}\); subtracting
the indicated central dimensions gives (D.51.7).

\emph{The exact common action and its central lift.}
By (D.51.2), the image of \(\mathfrak R_{50}\)
in \(\mathcal T_-\) is exactly \(Z/Z_0=Z_-\).
Products of the original source kernels therefore
give the maximal common quotient
\(\mathcal T_2/Z\).
Its action is also determined, not presumed.
The image of \(K_3\) in \(\mathcal Z_h\) is
invariant under the fixed D.50 action, since
the latter is identity on \(G_3\).
Its maximal class-two exponent-four quotient
is characteristic. The kernel in the original
source is precisely the product of
\(\mathfrak R_{50}\) with
\(\overline{\Gamma_3(K_3)(K_3')^2K_3^4}\).
Thus the resulting quotient is
\(\mathcal T_2/Z\), and the actual D.50
action and conjugators descend to it.

Lift these two conjugator elements to
\(\mathcal T_-\). Their torsion conjugates
have the required orders; their cusp error
is in \(Z_-\). Left multiplication of the
conjugators by \(b,c\in Z_-\) changes that
error by exactly
\[
             (b,c)\longmapsto
                    (1-s)b+s(1-r)c .
 \tag{D.51.15}
\]
This uses the abelian normal kernel \(Z_-\);
conjugation on it depends only on \(G_3\).
There are no omitted quadratic corrections.
The image is the augmentation image for
\(\langle s,r\rangle=G_3\). Since
\((Z_-)_A=0\), it is all of \(Z_-\).
Choose a correction killing the error.
It preserves the complete common-quotient
conjugators and fixes the cusp exactly.

We check both factorization and generation.
The corrected map from \(\widehat B\) sends
\(K_3\) into a class-two exponent-four group,
so first factors from \(\mathcal T_2\).
Its induced map \(T\) on \(V\) is the fixed
D.50 map and is \(G_3\)-linear, since the
base action is identity. Therefore it preserves
\(W=[V,A]\), and its exterior map preserves
\(Z^A=Z_0\). The map consequently kills
the additional relation \(Z_0\) in the target,
and is an endomorphism of \(\mathcal T_-\).
Its action on the Frattini quotient of
\(D/Z_0\) is the unchanged invertible \(T\).
For a finite two-group an onto Frattini map
is onto: a proper image would be contained
in a maximal subgroup, whose index is two
and which contains the Frattini subgroup.
Thus the endomorphism is an automorphism.

The two actual component automorphisms agree
on \(\mathcal T_2/Z\), as do their specified
conjugators. The closed fibre product gives
exactly (D.51.5), with one conjugator pair
on that source. Its parameter projects to
the entire D.50 parameter. Since the latter
retains the full loop abelianization and
has source kernel inside \(\widehat F''\),
the new parameter has an actual
\(\widehat F'\) preimage. This preserves
the normalized hexagons and all the earlier
comparison maps.

The fibre-product projections give the first
two kernel formulas in (D.51.6). Finally an
element of \(\mathfrak R_{50}\) maps into
the kernel of \(\mathcal T_2\to\mathcal T_-\)
exactly when its image is in \(Z_0\).
Every element of \(Z_0\) has such a preimage
by (D.51.2). This proves the final exact
image in (D.51.6).

\emph{The exact finite moving lift.}
The program \texttt{d\_cm\_class2\_moving.py}
reads the frozen D.50 certificate without
modifying it. It reconstructs the literal
order-\(1080\) base, all \(181\) free graph
generators as words in \(s,r,r^{-1}\), and
their collected class-two coordinates.
For a closed path, appending a free generator
\(g_j^{\pm1}\) changes its \(\mathbb Z/4\)
coordinate and adds the commutators with
the higher-index odd coordinates. This is
the faithful normal form of D.50, so the
word checks take place in the actual quotient.

Here \(A=\langle z\rangle\) is central of
order three. The program obtains an actual
order-three lift by taking the fourth power
of a chosen \(z\)-word; its cube is checked
in the full class-two group. On the moving
abelian summand, \((1-z)^{-1}=2-z\) modulo
four. The recorded D.50 conjugator coordinates
therefore lift to actual commutators with
this order-three element. Their abelian
coordinates remain exactly the recorded
modulo-four reductions of the D.50
modulo-sixteen conjugators.

Their cusp error lies in \(Z\), verified
in an adapted \(A\)-basis consisting of
61 fixed vectors and 60 nontrivial
two-dimensional blocks. The program splits
the error by \(1+z+z^2\). On the moving
central part the inverse of \(1+z\) is \(z\).
Writing the literal \(z\)-word in \(s,r\)
telescopes \((1-z)u\) into the two columns
of (D.51.15), giving explicit central
corrections. Each central coordinate is
then realized by its actual free-generator
commutator word.

The resulting \(s\)- and \(r\)-conjugates
are checked to have exact orders two and
three. Direct collection of all
\(4\,159\,444\) letters in their cusp word
gives exactly the invariant component of
the original error. It vanishes in
\(\mathcal T_-\). Every column of the old
181-dimensional Frattini matrix is unchanged
and has full rank. Thus the finite verification
checks a compatible automorphism, not merely
a central equation without generation.

\emph{The separate invariant obstruction in the finite certificate.}

The program \texttt{d\_class2\_diagnostic.py}
independently reconstructs the same marked
graph and reads only the frozen D.50 matrix.
The associated torsion-free cover has
36 cusps of width 30. Euler characteristic
therefore gives genus 73 and compact homology
rank 146.

For two integral graph cycles \(a,b\), its
oriented ribbon intersection form is
\[
       \Omega(a,b)
          ={1\over2}\sum_{g\langle r\rangle}
                   (a_g b_{gr}-b_g a_{gr}).
 \tag{D.51.16}
\]
This formula follows by averaging the two
small transverse push-offs of each cycle
in the thickened graph. The trivalent
vertices contribute the displayed
determinant; the bivalent vertices contribute
zero. Each crossing is then counted twice.
The cycle condition makes the expression
independent of the representative of each
cyclically ordered \(r\)-vertex. The program
checks integrality in the entire integral
cycle basis and obtains rank 146 modulo two.

A multiplier-one peripheral automorphism
inducing identity on this marked base fixes
each cusp conjugacy class. To see the
resulting class-two constraint directly,
use a surface presentation with the single
boundary relation
\(\prod_i[a_i,b_i]\prod_jc_j=1\).
After taking its class-two image and
quotienting homology by the span of the
cusp classes, conjugating any \(c_j\)
contributes a wedge with \(c_j\) and hence
vanishes. The compact symplectic tensor is
therefore preserved. Equivalently its
homology matrix satisfies
\(T^{\mathsf t}\Omega T=\Omega\) modulo two.
This constraint holds for every choice of
conjugators with that same homology action.

The exact D.50 finite matrix violates this
identity: the defect has rank 72 and entry
\((0,6)\) equal to one. It cannot lift
peripherally to the full class-two source.
The moving lift above retains the same
homology action and succeeds in the quotient
specified in (D.51.4). This computation
addresses only the literal finite certificate;
the actual arithmetic \(G_3\) and its chosen
D.50 action are covered by the symbolic
lift and exact residual calculations above.
\end{proof}
